% Lösungen Einheit 3 von 3 – Punktsymmetrie, Drehung, Verschiebung (zusammenbau v0.1)
\einheitenkopf{Einheit 3 von 3 · Punktsymmetrie, Drehung, Verschiebung}

\erg{22}{a) verschoben \quad b) Bilddreieck $A'(5|2)$, $B'(7|2)$, $C'(6|4)$ \quad c) $4$, $2$ – 4 Kästchen nach rechts, 2 nach unten \quad d) Bilddreieck $A'(3|3)$, $B'(1|3)$, $C'(3|1)$ \quad e) Bilddreieck $A'(-1|1)$, $B'(-1|4)$, $C'(-3|1)$ \quad f) $P'(7|4)$ \quad g) Bilddreieck $A'(7|3)$, $B'(5|2)$, $C'(6|1)$ \quad h) ja \quad i) Verschiebung nach rechts; nächstes Dreieck $(6|1)$, $(8|1)$, $(7|3)$, danach $(9|1)$, $(11|1)$, $(10|3)$ \quad j) Verschiebung (entlang beider Seitenrichtungen) \quad k) Drehung um Z (halbe Drehung, Punktspiegelung); beim Drehen bleiben alle Längen und Winkel gleich, nur die Lage ändert sich – also ist A'B' so lang wie AB.}
\erg{23}{a) $120^\circ$}
\erg{24}{a) $A'B' = 4{,}2$\,cm, Winkel bei $A' = 50^\circ$}
\erg{25}{a) Sie hat an einer Achse gespiegelt statt gedreht; bei der halben Drehung liegt jeder Bildpunkt auf der Geraden durch Z, gleich weit auf der anderen Seite. \quad b) Jeder Punkt wandert gleich weit in dieselbe Richtung; die Figur wird nur bewegt, nicht verzerrt. \quad c) $4$ nach rechts, 1 nach oben \quad d) $2 + 30 = 32$\,m}
